\documentclass[aspectratio=169,12pt]{beamer}
\usepackage{fontspec}
\usepackage[portuguese]{babel}
\usepackage{xcolor,listings,tikz,booktabs,amsmath,amssymb,graphicx,multicol,array,colortbl}
\usetikzlibrary{arrows.meta,positioning,shapes.geometric,calc}

\definecolor{navy}{HTML}{0D1B3E}
\definecolor{gold}{HTML}{F5A623}
\definecolor{qgreen}{HTML}{1E8B4C}
\definecolor{qblue}{HTML}{1A6EAE}
\definecolor{codebg}{HTML}{EEF3F8}

\usetheme{default}
\useinnertheme{circles}
\setbeamertemplate{navigation symbols}{}
\setbeamertemplate{footline}{%
  \hfill{\color{gray!60}\scriptsize\insertframenumber/\inserttotalframenumber}\hspace{.6em}\vspace{.4em}}

\setbeamercolor{frametitle}{bg=navy,fg=white}
\setbeamercolor{structure}{fg=navy}
\setbeamercolor{alerted text}{fg=gold}
\setbeamercolor{item}{fg=gold}
\setbeamercolor{subitem}{fg=navy}
\setbeamercolor{block title}{bg=qblue,fg=white}
\setbeamercolor{block body}{bg=qblue!8,fg=black}
\setbeamercolor{block title alerted}{bg=gold,fg=navy}
\setbeamercolor{block body alerted}{bg=gold!15,fg=navy}
\setbeamercolor{block title example}{bg=qgreen,fg=white}
\setbeamercolor{block body example}{bg=qgreen!10,fg=black}
\setbeamerfont{frametitle}{size=\large,series=\bfseries}
\setbeamerfont{block title}{series=\bfseries}

\setbeamertemplate{title page}{%
  \begin{tikzpicture}[remember picture,overlay]
    \fill[navy](current page.north west)rectangle(current page.south east);
    \fill[gold]([yshift=.7cm]current page.south west)rectangle(current page.south east);
    \node[anchor=south,yshift=.73cm,navy,font=\tiny\bfseries]at(current page.south)
      {INTENSIV\~{A}O QUANT AI \textbullet{} IMPACT UFSCAR};
  \end{tikzpicture}
  \vspace{.65cm}
  \begin{center}
    {\color{gold}\scriptsize\bfseries INTENSIV\~{A}O QUANT AI \textbullet{} IMPACT UFSCAR}\\[.4cm]
    {\color{white}\LARGE\bfseries\inserttitle}\\[.25cm]
    {\color{gold!85}\normalsize\insertsubtitle}\\[.5cm]
    {\color{white!70}\small\insertauthor}\\[.1cm]
    {\color{white!50}\footnotesize\insertdate}
  \end{center}}

\lstdefinestyle{python}{language=Python,
  basicstyle=\ttfamily\scriptsize\color{qblue},
  keywordstyle=\color{navy}\bfseries,
  stringstyle=\color{qgreen!80!black},
  commentstyle=\color{gray!70}\itshape,
  backgroundcolor=\color{codebg},
  frame=l,framerule=2pt,rulecolor=\color{gold},
  xleftmargin=1em,breaklines=true,showstringspaces=false,tabsize=4,
  extendedchars=false,inputencoding=ascii}
\lstset{style=python}

\author{Felipe Maldonado\\{\scriptsize VP Diretoria Quant~\textbullet{}~Impact UFSCAR}}
\date{Intensiv\~{a}o Quant AI 2025}
\title{Aula 01 --- Kickoff}
\subtitle{Do Zero ao Portfólio Quant em 10 Semanas}

\begin{document}

\begin{frame}[plain]
\titlepage
\end{frame}

\begin{frame}[fragile]{Agenda}

\begin{multicols}{2}
\tableofcontents
\end{multicols}
\end{frame}

\section{Motivação}
\begin{frame}[fragile]{Por que Gestão Quantitativa?}
\begin{columns}[T]
  \column{0.5\textwidth}
\begin{itemize}
  \item Retornos consistentes baseados em \textbf{processo}, não intuição
  \item Escala: uma estratégia testa \alert{milhares} de ativos simultaneamente
  \item Reprodutibilidade: resultados auditáveis e rastreáveis
  \item Vantagem competitiva: explora anomalias antes do mercado precificar
\end{itemize}
  \column{0.47\textwidth}
\begin{block}{Mercado Brasileiro}
\begin{itemize}
  \item B3: \textasciitilde 450 ações listadas
  \item IBOVESPA: \textasciitilde 90 ativos, revisado a cada 4 meses
  \item Fundos quant no Brasil cresceram $>$200\% em 5 anos
  \item Itaú, Verde, Giant Steps, Kadima, Ibiuna
\end{itemize}
\end{block}
\end{columns}
\end{frame}

\section{Universos de Investimento}
\begin{frame}[fragile]{Universos de Investimento}

\begin{center}
\renewcommand{\arraystretch}{1.3}
\begin{tabular}{>{\bfseries\color{navy}}l p{2.8cm} p{2.8cm} p{3cm}}
\rowcolor{navy}\color{white}\textbf{Classe} & \color{white}\textbf{Exemplos} &
  \color{white}\textbf{Retorno} & \color{white}\textbf{Risco principal}\\
\rowcolor{qblue!10}Renda Fixa     & Tesouro, CDB, LCI & Juros + spread    & Crédito / duration\\
Renda Variável  & Ações, FIIs, BDRs & Dividendo + capital& Volatilidade\\
\rowcolor{qblue!10}Derivativos    & Futuros, Opções   & Assimétrico       & Alavancagem\\
Alternativos    & FIPs, Crypto      & Alta variância    & Liquidez\\
\end{tabular}
\end{center}
\end{frame}

\begin{frame}[fragile]{Renda Fixa — Fundamentos}
\begin{columns}[T]
  \column{0.5\textwidth}
\begin{itemize}
  \item \textbf{Tesouro Selic} — pós-fixado, risco soberano mínimo
  \item \textbf{Tesouro IPCA+} — proteção contra inflação (duration longa)
  \item \textbf{CDB / LCI / LCA} — risco bancário, cobertura FGC até R\$250k
  \item \textbf{Debêntures} — spread de crédito, menor liquidez
\end{itemize}
  \column{0.47\textwidth}
\begin{block}{Conceitos-chave}
\begin{itemize}
  \item \textbf{Duration}: sensibilidade ao juro; $\Delta P \approx -D\,\Delta y$
  \item \textbf{Spread}: prêmio sobre a taxa livre de risco
  \item \textbf{Rating}: probabilidade implícita de default
\end{itemize}
\end{block}
\end{columns}
\end{frame}

\begin{frame}[fragile]{Renda Variável — Ações e Derivados}
\begin{columns}[T]
  \column{0.5\textwidth}
\begin{itemize}
  \item \textbf{Ações ordinárias (ON)}: direito a voto + dividendos
  \item \textbf{Ações preferenciais (PN)}: prioridade no dividendo
  \item \textbf{FIIs}: imóveis via bolsa, distribuição mensal obrigatória
  \item \textbf{BDRs}: recibos de ações estrangeiras negociadas na B3
\end{itemize}
  \column{0.47\textwidth}
\begin{block}{Componentes do retorno}
\begin{itemize}
  \item $r_t = \dfrac{P_t - P_{t-1} + D_t}{P_{t-1}}$
  \item $D_t$: dividendos e juros sobre capital próprio
  \item Retorno total ex-dividendos $\neq$ retorno total
\end{itemize}
\end{block}
\end{columns}
\end{frame}

\begin{frame}[fragile]{Derivativos e Mercados Alternativos}
\begin{columns}[T]
  \column{0.5\textwidth}
\begin{itemize}
  \item \textbf{Futuros}: obrigação de comprar/vender no vencimento (IBOV, DI, Dólar)
  \item \textbf{Opções}: direito (não obrigação); Call e Put
  \item \textbf{Swaps}: troca de fluxos de caixa (DI $\times$ IPCA, por ex.)
  \item \textbf{Hedge vs especulação}: mesmo instrumento, propósitos opostos
\end{itemize}
  \column{0.47\textwidth}
\begin{block}{Alternativos}
\begin{itemize}
  \item FIPs: private equity ilíquido, horizonte $>$5 anos
  \item Commodities: ouro, petróleo, agro — diversificação real
  \item Cripto: alta vol, correlação baixa com RV tradicional
\end{itemize}
\end{block}
\end{columns}
\end{frame}

\section{Teoria Financeira Clássica}
\begin{frame}[fragile]{Teoria Moderna do Portfólio — Markowitz (1952)}
\begin{columns}[T]
  \column{0.5\textwidth}
\textbf{Problema de otimização:}
\[
\min_{\mathbf{w}}\;\mathbf{w}^\top\boldsymbol{\Sigma}\mathbf{w}
\quad\text{s.t.}\quad
\mathbf{w}^\top\boldsymbol{\mu}\ge r^*,\;
\mathbf{w}^\top\mathbf{1}=1
\]
\vspace{.3em}
\begin{itemize}
  \item $\mathbf{w}$: vetor de pesos
  \item $\boldsymbol{\Sigma}$: matriz de covariâncias $n\times n$
  \item $\boldsymbol{\mu}$: vetor de retornos esperados
  \item Resultado: \alert{fronteira eficiente}
\end{itemize}
  \column{0.47\textwidth}
\begin{block}{Insight central}
Diversificação reduz risco sem necessariamente reduzir retorno. O risco de um portfólio depende das \textbf{covariâncias} entre ativos, não só das variâncias individuais.
\end{block}
\vspace{.3em}
\begin{alertblock}{Limitação prática}
Markowitz é extremamente sensível à estimação de $\boldsymbol{\mu}$ — pequenos erros geram pesos concentrados e instáveis (Michaud, 1989).
\end{alertblock}
\end{columns}
\end{frame}

\begin{frame}[fragile]{Hipótese de Mercado Eficiente — Fama (1970)}

\begin{columns}[T]
  \column{0.32\textwidth}
  \begin{block}{Forma Fraca}
    Preços refletem todo o histórico de preços e volumes.\\[.3em]
    \textbf{Implicação}: análise técnica não gera alpha.
  \end{block}
  \column{0.32\textwidth}
  \begin{block}{Forma Semiforte}
    + informações públicas (relatórios, dividendos).\\[.3em]
    \textbf{Implicação}: análise fundamentalista não gera alpha.
  \end{block}
  \column{0.32\textwidth}
  \begin{alertblock}{Forma Forte}
    + informações privadas (insider).\\[.3em]
    \textbf{Implicação}: ninguém bate o mercado consistentemente.
  \end{alertblock}
\end{columns}
\vspace{.4em}
\begin{exampleblock}{Por que estudamos anomalias então?}
  Evidências empíricas contra a HME: anomalias persistentes sugerem que mercados são
  \textit{aproximadamente} eficientes — e as ineficiências são nossa oportunidade.
\end{exampleblock}
\end{frame}

\section{Anomalias e Factor Investing}
\begin{frame}[fragile]{Anomalias de Mercado}
\begin{columns}[T]
  \column{0.5\textwidth}
\begin{itemize}
  \item \textbf{Size} (Banz, 1981): small caps superam large caps em risco ajustado
  \item \textbf{Value} (Stattman, 1980): alto P/VPA $\Rightarrow$ baixo retorno futuro
  \item \textbf{Momentum} (Jegadeesh \& Titman, 1993): vencedores continuam vencendo
  \item \textbf{Low Volatility} (Black, 1972): ativos de baixa vol superam CAPM
  \item \textbf{Quality} (Novy-Marx, 2013): empresas rentáveis superam
\end{itemize}
  \column{0.47\textwidth}
\begin{block}{Persistência no Brasil}
\begin{itemize}
  \item Mussa et al.\ (2012): momentum e value significativos no IBOVESPA
  \item Fama-French 3F aplicado ao Brasil por Málaga \& Securato (2004)
  \item Prêmios menores que EUA, mas estatisticamente presentes
\end{itemize}
\end{block}
\vspace{.3em}
\begin{alertblock}{Explicações}
\begin{itemize}
  \item Risco não capturado pelo CAPM
  \item Vieses comportamentais dos investidores
  \item Fricções de mercado (liquidez, custos)
\end{itemize}
\end{alertblock}
\end{columns}
\end{frame}

\begin{frame}[fragile]{Factor Investing — Fama \& French (1993)}
\begin{columns}[T]
  \column{0.5\textwidth}
\textbf{CAPM} (Sharpe, 1964):
\[ r_i - r_f = \alpha + \beta_{\text{MKT}}\,(r_m-r_f) + \varepsilon \]

\textbf{Fama-French 3F} (1993):
\[ r_i - r_f = \alpha + \beta_M\,\text{MKT} + \beta_S\,\text{SMB} + \beta_V\,\text{HML} + \varepsilon \]

\textbf{Fama-French 5F} (2015):
\[ +\;\beta_R\,\text{RMW} + \beta_I\,\text{CMA} \]

\vspace{.2em}
\begin{itemize}
  \item SMB: \textit{Small Minus Big}
  \item HML: \textit{High Minus Low} (book-to-market)
  \item RMW: \textit{Robust Minus Weak} (rentabilidade)
  \item CMA: \textit{Conservative Minus Aggressive} (investimento)
\end{itemize}
  \column{0.47\textwidth}
\begin{block}{Nossa estratégia}
\begin{itemize}
  \item Foco em \textbf{Momentum} — não coberto pelo modelo 5F original
  \item Carhart (1997) adicionou MOM como 4º fator
  \item Momentum é o fator com maior Sharpe histórico (Asness et al., 2013)
\end{itemize}
\end{block}
\vspace{.3em}
\begin{exampleblock}{Smart Beta}
\begin{itemize}
  \item ETFs que replicam fatores sistematicamente
  \item Ex: BOVA11 (mercado), SMALL11 (size)
  \item Fundos quant = seleção dinâmica de ativos
\end{itemize}
\end{exampleblock}
\end{columns}
\end{frame}

\section{Momentum e Behavioral Finance}
\begin{frame}[fragile]{O Fator Momentum — Jegadeesh \& Titman (1993)}
\begin{columns}[T]
  \column{0.5\textwidth}
\textbf{Definição do sinal 12-1:}
\[
\text{Sinal}_t = \sum_{k=2}^{12} r_{t-k}
\]
equivalente a:
\begin{lstlisting}
ret.shift(2).rolling(11).sum()
\end{lstlisting}
\vspace{.3em}
\begin{itemize}
  \item Janela de 12 meses; exclui mês mais recente (reversão de curto prazo)
  \item Ranking cross-sectional: compra top decil, vende bottom decil
  \item Horizonte de holding: 3--12 meses
\end{itemize}
  \column{0.47\textwidth}
\begin{block}{Evidências Empíricas}
\begin{itemize}
  \item EUA (1927--1989): prêmio de 3--12\% a.a.\ (J\&T, 1993)
  \item Internacional: 40 países (Rouwenhorst, 1998)
  \item Brasil: significativo em períodos $>$6 meses (Mussa et al., 2012)
  \item Ativos: ações, moedas, commodities, bonds (Asness et al., 2013)
\end{itemize}
\end{block}
\vspace{.3em}
\begin{alertblock}{Momentum Crash}
\begin{itemize}
  \item Crises de liquidez revertem momentum abruptamente
  \item Controle: ajuste por volatilidade (Moskowitz et al., 2012)
\end{itemize}
\end{alertblock}
\end{columns}
\end{frame}

\begin{frame}[fragile]{Finanças Comportamentais}
\begin{columns}[T]
  \column{0.5\textwidth}
\begin{itemize}
  \item \textbf{Herding}: investidores seguem a maioria, amplificando tendências
  \item \textbf{Underreaction}: mercado demora a incorporar novas informações
  \item \textbf{Overreaction}: eventual correção exagerada (reversão de longo prazo)
  \item \textbf{Disposition Effect}: vende ganhador cedo, segura perdedor
\end{itemize}
  \column{0.47\textwidth}
\begin{block}{Hong \& Stein (1999)}

\textit{A unified theory of underreaction, momentum trading and overreaction.}\\[.3em]
Dois tipos de agentes:
\begin{itemize}
  \item \textbf{News watchers}: agem em fundamentos, propagação lenta
  \item \textbf{Momentum traders}: agem em preço passado
\end{itemize}
Interação gera momentum no curto prazo e reversão no longo prazo.
\end{block}
\vspace{.2em}
\begin{exampleblock}{Kahneman (2011)}
Sistema 1 (rápido, intuitivo) vs Sistema 2 (lento, racional). Vieses cognitivos são previsíveis e exploráveis.
\end{exampleblock}
\end{columns}
\end{frame}

\section{O Desafio Itaú}
\begin{frame}[fragile]{O Desafio Itaú Asset Management}
\begin{columns}[T]
  \column{0.5\textwidth}
\begin{block}{O que é?}
\begin{itemize}
  \item Competição nacional de gestão quantitativa de portfólios
  \item Promovida pela Itaú Asset Management
  \item Foco em estratégias sistemáticas com dados reais
  \item Participantes: universidades e grupos de pesquisa
\end{itemize}
\end{block}
\vspace{.3em}
\begin{block}{Critérios de Avaliação}
\begin{itemize}
  \item \textbf{Rigor metodológico}: backtest sem viés, walk-forward
  \item \textbf{Fundamentação econômica}: por que o fator funciona?
  \item \textbf{Qualidade do código}: reprodutível, documentado
  \item \textbf{Comunicação}: clareza, honestidade sobre limitações
\end{itemize}
\end{block}
  \column{0.47\textwidth}
\begin{block}{Entregáveis}
\begin{itemize}
  \item Relatório técnico (PDF)
  \item Código reprodutível (GitHub)
  \item Apresentação oral de 10 min + 5 min de perguntas
  \item Pipeline completo: dados $\to$ sinal $\to$ portfólio $\to$ backtest $\to$ validação
\end{itemize}
\end{block}
\vspace{.3em}
\begin{exampleblock}{Nossa preparação}
\begin{itemize}
  \item 10 semanas de treinamento intensivo
  \item Pipeline completo em Python com dados reais do IBOVESPA
  \item Sinal de momentum academicamente fundamentado
  \item Backtest rigoroso com DSR e walk-forward
\end{itemize}
\end{exampleblock}
\end{columns}
\end{frame}

\section{Nossa Abordagem}
\begin{frame}[fragile]{Pipeline Técnico do Intensivão}

\begin{center}
\begin{tikzpicture}[node distance=.7cm and .5cm,
  box/.style={draw=navy,fill=navy!8,rounded corners=3pt,text width=1.7cm,align=center,
              font=\scriptsize\bfseries,inner sep=4pt},
  arr/.style={-Stealth,navy,thick}]
  \node[box](d){Dados\\{\tiny yfinance}};
  \node[box,right=of d](e){EDA\\{\tiny análise}};
  \node[box,right=of e](s){Sinal v1\\{\tiny 12-1}};
  \node[box,right=of s](p){Portfólio\\{\tiny top 20\%}};
  \node[box,right=of p](b){Backtest\\{\tiny IS}};
  \node[box,right=of b](v){Validação\\{\tiny OOS/DSR}};
  \node[box,right=of v](r){Relatório\\{\tiny defesa}};

  \draw[arr](d)--(e); \draw[arr](e)--(s); \draw[arr](s)--(p);
  \draw[arr](p)--(b); \draw[arr](b)--(v); \draw[arr](v)--(r);

  \node[below=.4cm of s,font=\scriptsize,gold]{\textbf{Sinal v2}};
  \node[below=.4cm of p,font=\scriptsize,gold]{\textbf{Vol-weight}};
  \node[below=.4cm of v,font=\scriptsize,gold]{\textbf{Walk-forward}};
\end{tikzpicture}
\end{center}
\vspace{.2em}
\begin{columns}[T]
  \column{0.33\textwidth}
  \begin{block}{Dados (Aulas 2--3)}
    IBOVESPA via yfinance\\Limpeza, retornos mensais\\EDA: fat tails, correlações
  \end{block}
  \column{0.33\textwidth}
  \begin{block}{Estratégia (Aulas 4--7)}
    Sinal momentum 12-1\\Alocação: EW, vol-weight\\Backtest e métricas
  \end{block}
  \column{0.33\textwidth}
  \begin{block}{Validação (Aulas 8--10)}
    Walk-forward OOS\\DSR, análise de sensibilidade\\GenAI + relatório final
  \end{block}
\end{columns}
\end{frame}

\begin{frame}[fragile]{Roadmap --- 10 Aulas}

\begin{center}
\renewcommand{\arraystretch}{1.25}
\begin{tabular}{>{\bfseries\color{navy}}c >{\bfseries}l p{5.5cm} l}
\rowcolor{navy}\color{white}\# & \color{white}Tema & \color{white}Conteúdo principal & \color{white}Entregável\\
\rowcolor{gold!15}01 & Kickoff        & Teoria, mercados, pipeline, roadmap        & Este slide\\
02 & Dados          & yfinance, limpeza, retornos mensais       & 3 parquets\\
\rowcolor{gold!15}03 & EDA            & Distribuições, fat tails, correlações     & Gráficos\\
04 & Sinal v1       & Momentum 12-1, IC, ranking                & sinal\_v1.parquet\\
\rowcolor{gold!15}05 & Backtest v1    & Sharpe, MDD, equity curve                 & ret\_carteira.parquet\\
06 & Alocação       & Equal-weight, vol-weight, Markowitz       & pesos\_v2.parquet\\
\rowcolor{gold!15}07 & Sinal v2       & Vol-adjusted momentum, IC melhorado       & sinal\_v2.parquet\\
08 & Backtest Rig.  & Walk-forward, DSR, look-ahead bias        & ret\_oos.parquet\\
\rowcolor{gold!15}09 & GenAI          & API Anthropic, narrativas, pesquisa       & Relatório narrativo\\
10 & Relatório      & Painel final, sensibilidade, defesa       & PDF + GitHub\\
\end{tabular}
\end{center}
\end{frame}

\begin{frame}[fragile]{Referências Bibliográficas}

\begin{multicols}{2}
\tiny
\begin{itemize}
  \item \textbf{Markowitz, H.} (1952). Portfolio Selection. \textit{Journal of Finance}, 7(1), 77--91.
  \item \textbf{Sharpe, W.} (1964). Capital asset prices. \textit{Journal of Finance}, 19(3), 425--442.
  \item \textbf{Fama, E.} (1970). Efficient capital markets. \textit{Journal of Finance}, 25(2), 383--417.
  \item \textbf{Black, F.} (1972). Capital market equilibrium with restricted borrowing. \textit{Journal of Business}, 45(3).
  \item \textbf{Banz, R.} (1981). The relationship between return and market value. \textit{JFE}, 9(1), 3--18.
  \item \textbf{Stattman, D.} (1980). Book values and stock returns. \textit{Chicago MBA}, 4, 25--45.
  \item \textbf{Jegadeesh, N. \& Titman, S.} (1993). Returns to buying winners. \textit{Journal of Finance}, 48(1), 65--91.
  \item \textbf{Fama, E. \& French, K.} (1993). Common risk factors. \textit{JFE}, 33(1), 3--56.
  \item \textbf{Rouwenhorst, K.} (1998). International momentum strategies. \textit{Journal of Finance}, 53(1), 267--284.
  \item \textbf{Hong, H. \& Stein, J.} (1999). A unified theory of underreaction. \textit{Journal of Finance}, 54(6), 2143--2184.
  \item \textbf{Carhart, M.} (1997). On persistence in mutual fund performance. \textit{Journal of Finance}, 52(1), 57--82.
  \item \textbf{Kahneman, D.} (2011). \textit{Thinking, Fast and Slow}. Farrar, Straus and Giroux.
  \item \textbf{DeMiguel, V.\ et al.} (2009). Optimal versus naive diversification. \textit{RFS}, 22(5), 1915--1953.
  \item \textbf{Fama, E. \& French, K.} (2015). A five-factor asset pricing model. \textit{JFE}, 116(1), 1--22.
  \item \textbf{Moskowitz, T.\ et al.} (2012). Time series momentum. \textit{JFE}, 104(2), 228--250.
  \item \textbf{Asness, C.\ et al.} (2013). Value and momentum everywhere. \textit{Journal of Finance}, 68(3), 929--985.
  \item \textbf{Novy-Marx, R.} (2013). The other side of value. \textit{JFE}, 108(1), 1--28.
  \item \textbf{Mussa, A.\ et al.} (2012). Hipótese de mercados eficientes e finanças comportamentais. \textit{REGE}, 19(2).
  \item \textbf{Bailey, D. \& López de Prado, M.} (2014). The deflated Sharpe ratio. \textit{Journal of Portfolio Management}, 40(5).
  \item \textbf{López de Prado, M.} (2018). \textit{Advances in Financial Machine Learning}. Wiley.
\end{itemize}
\end{multicols}
\end{frame}

\end{document}
